% !TeX spellcheck = en_us
\documentclass[a4paper,12pt,fleqn]{article}

\usepackage[left=1.50cm, right=1.50cm, top=3cm, bottom=2cm, 
headheight= 2.5cm, headsep= 4pt, footskip=25pt]{geometry}
\usepackage{fancyhdr}
\usepackage{lmodern}
\usepackage[table]{xcolor}
\usepackage{amssymb,amstext,amsmath}
\usepackage{eso-pic}
\usepackage{booktabs}
\usepackage{bm}
\usepackage{setspace}
\usepackage{graphicx}
\usepackage[colorlinks]{hyperref}
\setlength{\parindent}{0mm}
\setlength{\parskip}{2mm}
\usepackage{enumitem}
\usepackage{titlesec}
\usepackage{gensymb}
\usepackage{pgfgantt}
\usepackage{lscape}
\usepackage[absolute]{textpos}
\usepackage{comment}
\usepackage{cite}
\usepackage{tikz}
\usepackage{pgfplots}
\usetikzlibrary{shapes.geometric}
\usepackage{array}
\newcolumntype{P}[1]{>{\centering\arraybackslash}p{#1}}
\usepackage{longtable}
\usepackage[font={small}]{caption}
\usepackage[most]{tcolorbox}
\usepackage{tikz}
\usepackage{stmaryrd}
\usepackage{wrapfig}
%\usepackage[none]{hyphenat}
\usepackage{ctable}
\usepackage{ragged2e}
\usepackage[official]{eurosym}
\newcommand*\circled[1]{\tikz[baseline=(char.base)]{
		\node[shape=circle,draw,inner sep=1pt] (char) {#1};}}
\usepackage[symbol]{footmisc}
\renewcommand{\thefootnote}{\fnsymbol{footnote}}
\renewcommand{\thempfootnote}{\fnsymbol{mpfootnote}}
\renewcommand*\footnoterule{}

\usepackage{xcolor}
\definecolor{docColor}{cmyk}{1.00,1.00,0, 0.40}  % 0.90\% of black
\color{docColor}

\newtcolorbox{myhlbox}[2][]{colback=docColor!05, colframe=docColor!20, width=1.0\linewidth,
	boxrule=0pt, left=1em, right=1em, top=1em, bottom=1em,
	colbacktitle=docColor!20, coltitle=docColor,enhanced, attach boxed title to top left={xshift=2mm, yshift=-1mm},
	title=#2, #1}

\hypersetup{
	colorlinks=true,
	citecolor=blue,
	filecolor=magenta,
	linkcolor=blue,
	urlcolor=blue,
	pdftitle={SeminarAnnouncements},
	%bookmarks=true,,
}

\graphicspath{{./figs/}}

\definecolor{dgreen}{rgb}{0.0, 0.5, 0.0}

\allowdisplaybreaks

\newcommand*\mystrut[1]{\vrule width0pt height0pt depth#1\relax}
\renewcommand{\thesection}{\arabic{section}.}
\renewcommand{\thesubsection}{\arabic{section}.\arabic{subsection}}
\renewcommand{\thesubsubsection}{\arabic{section}.\arabic{subsection}.\arabic{subsubsection}}
\titleformat{\section}{\sffamily\large\bfseries}{\thesection}{1.0em}{}
\titleformat{\subsection}{\sffamily\large\bfseries}{\thesubsection}{1.0em}{}
\titleformat{\subsubsection}{\sffamily\normalsize\bfseries}{\thesubsubsection}{1.0em}{}
\titlespacing\section{0pt}{0pt plus 2pt minus 2pt}{0pt plus 2pt minus 2pt}
\renewcommand{\baselinestretch}{1.10}

\newcommand\CircledImage[1]{%
	\tikz\path[fill overzoom image={#1}, line width=0mm]circle[radius=0.5\linewidth];%
}

\input{format/defs}

\fancypagestyle{titlepage}{
	\fancyhf{}
	\insertLXheader
}

\fancypagestyle{plain}{%
	\fancyhf{}
	\renewcommand{\headrulewidth}{0pt}
	\renewcommand{\footrulewidth}{0pt}
	\insertnormalpage
}

\pagestyle{plain}

\newcommand\TitleBackground{%
	\AddToShipoutPictureBG*{
		\begin{tikzpicture}[remember picture,overlay]
			\node[opacity=0.30, anchor=center, rotate=-90] 
			at (current page.center)
			{\includegraphics[height=1.5\paperheight,keepaspectratio]{/Users/kannan/Library/CloudStorage/Dropbox/Kannan/0-templates/00-logos/bg-img/lms/lms-bg-1-portrait.pdf}};
		\end{tikzpicture}
	}
}

\begin{document}
	\sffamily
	\thispagestyle{titlepage}
	\TitleBackground
	\vspace*{-2em}
	\begin{center}
		\huge \textbf{LMS Seminar}
	\end{center}
	\begin{center}
		\Large 
		\textbf{The geometric nature of residual stress in growing tissues} 
	\end{center}
	\begin{center}
		\Large Alexander Erlich \\ 
		{\large Laboratoire Interdisciplinaire de Physique \\
		Universit{\'e} Grenoble Alpes} \\
		\vspace*{0.5em}
		{\large \textbf{Date, time, and venue} \\
		September 24, 2026 (2 -- 3 pm), Amphithéâtre Pole Mecanique}
	\end{center}

	\vspace*{-1em}
	
	\begin{myhlbox}{\textbf{Abstract}}
		This presentation investigates how organisms reach a specific size, focusing on the role of residual stress, which remains in tissues even after external forces are removed. Currently, the respective role of mechanics and biochemical processes in determining organ size is an open question, and it is unknown how stress is actively built and maintained in tissues during growth.

		The central hypothesis of this talk suggests growth incompatibility as a missing link between tissue growth, size, and residual stress. Growth incompatibility represents the challenge of fitting parts of a grown tissue without voids or overlaps, acting as the geometric "seed" of residual stress. We use the Ricci curvature of the growth metric tensor to approximate incompatibility at both tissue and cell levels.

		We propose a theoretical framework where growth incompatibility acts as a geometric regulator for size termination during tissue development. Inspired by vertex models of morphogenesis, initially, we explore the hypothesis that the Ricci scalar curvature is prescribed in space and time [1]. Under this assumption, our model successfully reproduces specific experimental observations, such as the characteristic opening patterns observed after tissue cuts in Drosophila wing discs and multicellular spheroids, including the curvature of cut edges and consistent opening patterns following repeated cuts.

		Traditional biological growth models rely on the homeostatic (Eshelby-like) stress tensor to define an ideal target state. Any deviation from this state triggers growth and remodeling, aimed at restoring balance between mechanical forces and biological adaptation. Yet this homeostatic stress lacks clear biological interpretation and is often arbitrarily prescribed. To address this limitation, we shift our focus to growth incompatibility, removing the constraint of fixed Ricci curvature used earlier. Instead, we propose a formulation [2] that penalizes deviations from a target incompatibility state, analogous to the Einstein-Hilbert action in General Relativity. This provides a biologically meaningful and physically grounded approach that hints at a link between cellular mechanics and tissue-scale regulation of stress and size.
		
	\end{myhlbox}

\vspace*{1em}

	\begin{myhlbox}{\textbf{About the speaker}}
		\begin{wrapfigure}{r}{0.25\textwidth}
			\vspace*{-2em}
			\CircledImage{images/erlich}
		\end{wrapfigure}

		Alexander Erlich is an applied mathematician who investigates how mechanics shapes and regulates living biological systems. He is currently a CNRS researcher at the Laboratoire Interdisciplinaire de Physique in Grenoble, France, where he focuses on modelling growth in living elastic tissues. He is especially interested in how feedback laws between growth, mechanical stress, and chemical fields determine the final size of an organism. Alexander obtained his PhD in non-linear elasticity and growth from the University of Oxford, where he worked with Alain Goriely and Derek E. Moulton. He then did post- docs in fluid mechanics and transport at the University of Manchester, cell motility at LiPhy (Grenoble), and developmental biology at CENTURI (Marseille). His publications include topics such as growth-induced instabilities, placental transport, tissue folding, tumor spheroids, water lilies and Ammonites’ shells. He also collaborates with experimental biologists and physicists to test and validate his models.
		
	\end{myhlbox}
\end{document}